\section{Rates}
\setcounter{equation}{0}



Throughout this section $p$, $q$, $\sigma$, $B_X$, $c_U$, $C_U$, $\kappa_1$,
$\kappa_2$, $C_g$, $s$, $\pi$, $r$ and $\alpha$ are fixed as $n$ grows,
$J = J_n \to \infty$, $d_n = pq(2^{J_n}-1)$,
$2^{J_n} \le 2n$,
$\bn = \lceil\log n\rceil$ with $n \ge 3$, so that $\bn \ge 2$, $\Gcal_{n}$ is
the partition of Section~2.3 of the paper with $b = \bn$, the weights have
ratio $\varrho \le \bar\varrho$ for a fixed $\bar\varrho$, and $\pen = \lamG$
is as in (3.1).  The statements in $O_p$
are proved by exhibiting, for each $\epsilon > 0$, an event of probability at
least $1-\epsilon$ for $n$ large, built with the confidence levels
$\alpha = \alpha' = \alpha'' = \epsilon/4$ for the noise, the bias and the
levels, on which the error is at most a constant $M_\epsilon$ times the rate.
The constant $\Lambda = \kappa_2C_U + \gamma/2$ and the event $\ESig$ are those
of Section~S1.

% derivations/08-blocos.tex, as duas consequencias da Hipotese B.

\begin{lemma}[The penalty level and the regime]
\label{lem:regime}
In the setting above,
\begin{equation}
  \label{eq:order-lambda}
  (\lamGone)^{2} \;\le\; C_{\pen}\,\frac{\sigma^{2}\bn}{n},
  \qquad
  C_{\pen} = 16\bigl\{2B_X^{2}C_U + \Lambda\bigl(1 + \tfrac12\log(2pq/\alpha)\bigr)\bigr\},
\end{equation}
and, on the event $\Acal = \{\smax^{2} \le 2B_X^{2}C_U\}\cap\ESig$,
$\lamG \ge \pen_{w}$, so that $\Ecal \subseteq \Tw$ for the event $\Ecal$ of
Lemma~\ref{lem:calib-blocks}(i), $\gtil > \gamma/2$, $\Amat$ has rank $p$ and
$\lmaxx(\Gramhat) < \Lambda$. Moreover $\Prob(\Acal) \to 1$ whenever
$pq2^{J_n}\log(pq2^{J_n})/n \to 0$.
\end{lemma}

\begin{proof}
By $(x+y)^{2} \le 2x^{2}+2y^{2}$, $(\lamGone)^{2} = 4\max_{G}\lamzbar{G}^{2} \le
(8\sigma^{2}/n)\{2B_X^{2}C_U b_{\max} + 2\Lambda\log(|\Gcal_{n}|/\alpha)\}$, with
$b_{\max} = \max_{G}|G|$. By Proposition~\ref{prop:balanced},
$b_{\max} \le 2\bn - 1 < 2\bn$ and $|\Gcal_{n}| \le pq(1 + 2^{J_n}/\bn) \le
pq(1+n) \le 2pqn$, because $2^{J_n} \le 2n$ and $\bn \ge 2$; hence
$\log(|\Gcal_{n}|/\alpha) \le \log(2pq/\alpha) + \log n \le \log(2pq/\alpha) +
\bn$, and $(\lamGone)^{2} \le (16\sigma^{2}/n)\{(2B_X^{2}C_U + \Lambda)\bn +
\Lambda\log(2pq/\alpha)\} \le C_{\pen}\sigma^{2}\bn/n$, using $\bn \ge 2$. On
$\Acal$, Lemma~\ref{lem:calib-blocks}(ii) gives $\lamz{G} \le \lamzbar{G}$ for
every $G$, so that $\lamG \ge \pen_{w}$ and, by the same lemma,
$\Ecal \subseteq \Tw$. On $\ESig$, $\lmin(\Gramhat) > \gamma/2$ and
$\lmaxx(\Gramhat) < \Lambda$; \eqref{eq:schur} gives $\gtil > \gamma/2$, and
$\lmin(\Amat\T\Amat/n) \ge \lmin(\Gramhat) > 0$ gives the rank. Finally
$\Prob(\smax^{2} \le 2B_X^{2}C_U) \to 1$ by Lemma~\ref{lem:columns}, which asks
$2^{J}\log d/n \to 0$, and $\Prob(\ESig) \to 1$ when
$pq2^{J_n}\log(pq2^{J_n})/n \to 0$ (Section~S1); the second condition implies
the first.
\end{proof}



% derivations/08-blocos.tex, prova do Lema 15 (numeracao global).

\subsection*{Proof of Lemma 3}

Fix a block and write $\theta_{\lev} = \theta^{*}_{\cv\md,\lev\cdot} \in
\R^{2^{\lev}}$, set to zero for $\lev \ge J$, which only decreases the risk;
write $\Gcal_{\lev}$ for the chunks of a fine level $\lev$, those with
$2^{\lev} \ge b$, and $r_{\lev} = \sum_{G\in\Gcal_{\lev}}
\min(\norm{\theta_{G}}_{2}^{2},\eta)$. The coarse chunk, if there is one,
contributes at most $\eta$, so the block contributes at most
$\eta + \sum_{\lev}r_{\lev}$, the sum running over the fine levels. The first
bound on $r_{\lev}$ is the count: each chunk of $\Gcal_{\lev}$ has at least $b$
of the $2^{\lev}$ coefficients of the level, so that
$r_{\lev} \le |\Gcal_{\lev}|\eta \le 2^{\lev}\eta/b$.

\emph{The case $\pi \le 2$.} Write $a = \pi(s+1/2) > 1$. For $u \ge 0$,
$\min(u,\eta) \le \eta^{1-\pi/2}u^{\pi/2}$, and $\norm{x}_{2} \le \norm{x}_{\pi}$
for $\pi \le 2$; hence the second bound,
\[
  r_{\lev} \;\le\; \eta^{1-\pi/2}\sum_{G\in\Gcal_{\lev}}\norm{\theta_{G}}_{\pi}^{\pi}
  \;\le\; \eta^{1-\pi/2}\norm{\theta_{\lev}}_{\pi}^{\pi}
  \;\le\; \eta^{1-\pi/2}C_g^{\pi}2^{-\lev(a-1)},
\]
by Assumption~5, because the chunks of $\Gcal_{\lev}$ are disjoint and
contained in the level. Define $x \in \R$ by
$2^{xa} = bC_g^{\pi}\eta^{-\pi/2}$, so that the two bounds agree at
$\lev = x$, with common value
$h = 2^{x}\eta/b = \eta^{1-\pi/2}C_g^{\pi}2^{-x(a-1)}$. Summing the first bound
over the levels $\lev \le x$ gives at most $2^{\lfloor x\rfloor+1}\eta/b \le 2h$
(an empty sum if $x < 0$), and summing the second over the levels $\lev > x$
gives at most $h\sum_{i\ge0}2^{-i(a-1)} = h/(1-2^{1-a})$, since the first
integer $\lev > x$ has $2^{-\lev(a-1)} \le 2^{-x(a-1)}$. The block therefore
contributes at most $\eta + \{2 + (1-2^{1-a})^{-1}\}h$. Since $1/a = \tau/\pi$,
$2^{x} = b^{\tau/\pi}C_g^{\tau}\eta^{-\tau/2}$ and
$h = b^{\tau/\pi-1}C_g^{\tau}\eta^{1-\tau/2}$, and, with
$1-\tau/2 = 2s/(2s+1)$,
\[
  h \;=\; C_g^{\tau}\Bigl(\frac{\eta}{b}\Bigr)^{2s/(2s+1)}
  b^{\tau/\pi - 1 + 2s/(2s+1)}
  \;=\; C_g^{\tau}\Bigl(\frac{\eta}{b}\Bigr)^{2s/(2s+1)}
  b^{\tau(1/\pi-1/2)},
  \qquad
  \tau\Bigl(\frac{1}{\pi}-\frac12\Bigr) = \frac{2/\pi-1}{2s+1}.
\]

\emph{The case $\pi \ge 2$.} Hölder's inequality in $\R^{2^{\lev}}$ gives
$\norm{\theta_{\lev}}_{2} \le 2^{\lev(1/2-1/\pi)}\norm{\theta_{\lev}}_{\pi} \le
C_g2^{-\lev s}$, with $1/\pi = 0$ if $\pi = \infty$, and the second bound is
$r_{\lev} \le \norm{\theta_{\lev}}_{2}^{2} \le C_g^{2}2^{-2\lev s}$. With
$2^{x(2s+1)} = C_g^{2}b/\eta$ and $h = 2^{x}\eta/b = C_g^{2}2^{-2xs}$, the same
computation gives the contribution $\eta + \{2 + (1-2^{-2s})^{-1}\}h$, and
$h = (C_g^{2}b/\eta)^{1/(2s+1)}\eta/b = C_g^{\tau}(\eta/b)^{2s/(2s+1)}$, because
$2/(2s+1) = \tau$.

Summing over the $pq$ blocks gives the statement; at $\pi = 2$,
$1-2^{1-a} = 1-2^{-2s}$ and the two constants agree. \qed


We also record that, on the event of Proposition 2(ii) with $t = \gamma/2$,
\begin{equation}
  \label{eq:lmax}
  \normn{\Bmat\delta}^2 \le \lmaxx(\Gramhat)\norm{\delta}_2^2
  \quad\text{and}\quad \lmaxx(\Gramhat) \le \Lambda,
\end{equation}
the first because $\Bmat\T\Bmat/n$ is a principal submatrix of $\Gramhat$ and
the second by Weyl's inequality and Proposition 2(i).

% derivations/08-blocos.tex, a contagem (eq. contagem-lenta) e a prova do
% Corolario 14 (numeracao global); D47(c), D49(b).

\begin{lemma}[The norm of the oracle over chunks]
\label{lem:count}
Under Assumptions~4 and 5, let $\Gcal$ be the partition of Section~2.3 of the
paper with $b \ge 2$, and let $\vartheta_{\pi} = \min(1-1/\pi,\,1/2)$, with
$\vartheta_{\infty} = 1/2$. Then, for every $J$,
\begin{equation}
  \label{eq:count}
  \nGone{\thetastar} \;\le\; pqC_g\Bigl\{\frac{1}{1-2^{-\seff}}
  + b^{-\vartheta_{\pi}}\sum_{\lev=\jst+1}^{J-1}2^{\lev(1/2-s)}\Bigr\}.
\end{equation}
The bound for each fine level is attained, by a level whose energy is spread
evenly when $\pi \ge 2$, and by a level with one large coefficient in each
chunk when $\pi \le 2$.
\end{lemma}

\begin{proof}
Fix a block and write $\theta_{\lev}$ for its coefficients of level $\lev$.
By Assumption~5, $\norm{\theta_{\lev}}_{2} \le
2^{\lev(1/2-1/\pi)_{+}}\norm{\theta_{\lev}}_{\pi} \le C_g2^{-\lev\seff}$, by
Hölder's inequality in $\R^{2^{\lev}}$ if $\pi \ge 2$ and by
$\norm{\cdot}_{2} \le \norm{\cdot}_{\pi}$ if $\pi \le 2$. The coarse chunk
contributes at most $\sum_{\lev\le\jst}\norm{\theta_{\lev}}_{2} \le
C_g(1-2^{-\seff})^{-1}$. At a fine level $\lev$ the $k_{\lev} \le 2^{\lev}/b$
chunks $G \in \Gcal_{\lev}$ partition the level. If $\pi \ge 2$, by the
Cauchy-Schwarz inequality,
\[
  \sum_{G\in\Gcal_{\lev}}\norm{\theta_{G}}_{2}
  \;\le\; k_{\lev}^{1/2}\norm{\theta_{\lev}}_{2}
  \;\le\; \Bigl(\frac{2^{\lev}}{b}\Bigr)^{1/2}C_g2^{-\lev s}
  \;=\; C_g\,b^{-1/2}\,2^{\lev(1/2-s)} ;
\]
if $\pi \le 2$, by $\norm{\theta_{G}}_{2} \le \norm{\theta_{G}}_{\pi}$ and
Hölder's inequality over the chunks,
\[
  \sum_{G\in\Gcal_{\lev}}\norm{\theta_{G}}_{2}
  \;\le\; k_{\lev}^{1-1/\pi}\Bigl(\sum_{G\in\Gcal_{\lev}}
  \norm{\theta_{G}}_{\pi}^{\pi}\Bigr)^{1/\pi}
  = k_{\lev}^{1-1/\pi}\norm{\theta_{\lev}}_{\pi}
  \;\le\; \Bigl(\frac{2^{\lev}}{b}\Bigr)^{1-1/\pi}C_g2^{-\lev(s+1/2-1/\pi)}
  = C_g\,b^{-(1-1/\pi)}\,2^{\lev(1/2-s)} .
\]
Both are $C_gb^{-\vartheta_{\pi}}2^{\lev(1/2-s)}$, and the equality cases of
the two inequalities are the configurations of the statement. Summing over the
fine levels $\jst < \lev < J$ and over the $pq$ blocks gives \eqref{eq:count}.
\end{proof}

\subsection*{Proof of Proposition 3}

\emph{The nonasymptotic bound.} On $\Tw$, Theorem~2(i) with the comparator
$\bar\bbeta = \betastar$, so that $\bbar = \bb$ and
$\vbar = v = \thetahat-\thetastar$, gives
$\normn{\Btil v}^{2} \le 6\pen\nGw{\thetastar} + 4\normn{\bb}^{2}$, and the
first part of Theorem~2(iii), with $(x+y+z)^{2} \le 3(x^{2}+y^{2}+z^{2})$, gives
\begin{equation}
  \label{eq:slow-nonasymptotic}
  \normn{\widehat{f}-f}^{2} \;\le\; 18\pen\nGw{\thetastar}
  + 24\normn{\bb}^{2} + 3\normn{\PA\bveps}^{2}
  \qquad\text{on } \Tw .
\end{equation}
Steps 1 to 3 and 6 of the proof of Theorem~2, the only ones used, do not
involve $\Gramhat$. With $\pen \ge \pen_{w}$, calibrated on the covariates by
Lemma~\ref{lem:calib-blocks}(i), the bound holds with probability at least
$1-\alpha$ in any design in which $\Amat$ has rank $p$, including designs with
$d > n$, where $\gtil = 0$.

\emph{The penalty level of the theory.} Let
$\Acal^{+} = \{\smax^{2} \le 2B_X^{2}C_U\}\cap\{\lmaxx(\Gramhat) \le \Lambda\}$.
On it Lemma~\ref{lem:calib-blocks}(ii) gives $\lamz{G} \le \lamzbar{G}$ for
every $G$, hence $\lamG \ge \pen_{w}$ and $\Ecal \subseteq \Tw$. The first
event has probability tending to one by Lemma~\ref{lem:columns}, whose proof
uses only $|X_{\cv}| \le B_X$ and $p_{\bU} \le C_U$. For the second,
$\Prob\{\lmaxx(\Gramhat) \le \Lambda\} \ge
\Prob\{\opnorm{\Gramhat-\Gram} < \Lambda - \kappa_2C_U\}$ by Weyl's inequality,
since $\lmaxx(\Gram) \le \kappa_2C_U$, and the matrix Bernstein bound in the
proof of Proposition~2(ii), with $t = \Lambda - \kappa_2C_U$, sends it to one;
that bound uses only $\norm{\bZ_{i}}_{2}^{2} \le R_J$ and
$\lmaxx(\Gram) \le \kappa_2C_U$, the upper half of Proposition~2(i), and
neither $c_U$ nor $\kappa_1$, so that any $\Lambda > \kappa_2C_U$ serves. Both
probabilities tend to one when $pq2^{J_n}\log(pq2^{J_n})/n \to 0$. Fix
$\epsilon > 0$, take $\alpha = \alpha' = \alpha'' = \epsilon/4$ and intersect
$\Acal^{+}$, $\Ecal$, $\{\normn{\bb}^{2} \le \mathcal{B}_{J_n}^{2}/\alpha'\}$
and $\{\normn{\PA\bveps}^{2} \le \sigma^{2}p/(n\alpha'')\}$, the last two by
Markov's inequality, Lemma~2 and Theorem~2(iii); for $n$ large the
intersection has probability at least $1-\epsilon$, and on it
$\normn{\widehat{f}-f}^{2} \le 18\lamG\nGw{\thetastar} +
24\mathcal{B}_{J_n}^{2}/\alpha' + 3\sigma^{2}p/(n\alpha'')$, which is the first
display of Proposition~3. For the chain, Lemma~\ref{lem:calib-blocks}(iii)
gives $\lamG w_{G} \le \varrho\lamGone$ for every $G$, and summing with the
weights $\norm{\thetastar_{G}}_{2}$ gives
$\lamG\nGw{\thetastar} \le \varrho\lamGone\nGone{\thetastar}$, with
$\varrho \le \bar\varrho$; and $\norm{x}_{2} \le \norm{x}_{1}$ in each chunk
gives $\nGone{\thetastar} \le \norm{\thetastar}_{1}$. The order of $\lamGone$
is \eqref{eq:order-lambda}.

\emph{Part (i).} With $s > 1/2$ and $2^{\jst+1} \ge b$,
$\sum_{\lev>\jst}2^{\lev(1/2-s)} \le 2^{(\jst+1)(1/2-s)}/(1-2^{1/2-s}) \le
\bn^{1/2-s}/(1-2^{1/2-s})$, and Lemma~\ref{lem:count} gives
$\nGone{\thetastar} \le pqC_g\{(1-2^{-\seff})^{-1} +
\bn^{1/2-s-\vartheta_{\pi}}(1-2^{1/2-s})^{-1}\} = O(1)$; hence
$\lamG\nGw{\thetastar} = O\{(\bn/n)^{1/2}\} = O\{(\log n/n)^{1/2}\}$. With
$J_n \ge \log_2 n/(4\seff)$, $\mathcal{B}_{J_n}^{2} \asymp 2^{-2J_n\seff} \le
n^{-1/2}$, and $p/n$ is of smaller order. With $\seff > 1/4$ the choice
$J_n = \lceil\log_2 n/(4\seff)\rceil$ has $2^{J_n} < 2n^{1/(4\seff)}$ and
satisfies the conditions of the statement.

\emph{Part (ii).} With $s < 1/2$,
$\sum_{\jst<\lev<J_n}2^{\lev(1/2-s)} \le 2^{J_n(1/2-s)}/(2^{1/2-s}-1)$, so that
$\nGone{\thetastar} = O\{1 + \bn^{-\vartheta_{\pi}}2^{J_n(1/2-s)}\}$ and, since
$1/2 - \vartheta_{\pi} = (1/\pi-1/2)_{+}$,
\[
  \lamG\nGw{\thetastar} \;\lesssim\; \Bigl(\frac{\bn}{n}\Bigr)^{1/2}
  + n^{-1/2}\,\bn^{(1/\pi-1/2)_{+}}\,2^{J_n(1/2-s)} .
\]
Equating the second term with the bias $2^{-2J_n\seff}$ gives
$2^{J_n(1/2-s+2\seff)} \asymp n^{1/2}\bn^{-(1/\pi-1/2)_{+}}$, which is the
choice of the statement, with common value $2^{-2J_n\seff} \asymp
\{\bn^{(2/\pi-1)_{+}}/n\}^{2\seff/(1-2s+4\seff)}$. The term $(\bn/n)^{1/2}$ and
$p/n$ are of smaller order, because $2\seff/(1-2s+4\seff) < 1/2$ is equivalent
to $s < 1/2$. The conditions of the statement hold because $\seff > s/2$ gives
$1/(1-2s+4\seff) < 1$, and $2^{J_n} \ge \bn$ for $n$ large. When $\pi \ge 2$,
$\seff = s$, $(2/\pi-1)_{+} = 0$ and $1-2s+4s = 2s+1$. \qed

% derivations/08-blocos.tex, Corolario 14(iii) e (iv), a equacao mu-blocos e a
% prova dos dois itens (E1.14; E5e, D52).
\subsection*{Proposition 3 above the concentration ceiling}

The conditions $\seff > 1/4$ in Proposition~3(i) and $\seff > s/2$ in
Proposition~3(ii) keep the resolution of the balance below the ceiling
$2^{J_n} \asymp n/\log n$, under which $\smax^{2} \le 2B_X^{2}C_U$ and
$\lmaxx(\Gramhat) \le \Lambda$ hold with probability tending to one. The pairs
that they exclude, all with $\pi < 2$, and the boundary $s = 1/2$ call for
$2^{J_n}$ above that ceiling, where neither bound holds: the wavelets of the
finest level take values up to $2^{(J-1)/2}\norm{\psi}_{\infty}$ and have one
or two sample points in their support, so that $\smax^{2}$ and
$\lmaxx(\Gramhat)$ grow with $2^{J}/n$, and the deterministic level (3.1) is
no longer admissible. With $D = p + d$, $\xi = \Lambda - \kappa_2C_U > 0$,
$C_\psi$ as in Proposition~2(ii) and $\norm{\psi}_{\infty}$ as in
Lemma~\ref{lem:columns}, write
\begin{equation}
  \label{eq:mu-blocks}
  \mu^{\Gcal}_{J} = 1 + c_{\mu}\frac{2^{J}\log(2D)}{n},
  \qquad
  \lamGp = (\mu^{\Gcal}_{J_n})^{1/2}\lamG,
\end{equation}
with
\[
  c_{\psi} = \frac{7\norm{\psi}_{\infty}^{2}}{12C_U},
  \qquad
  c_{\mu} = \max\Bigl\{c_{\psi},\
  \frac{2pB_X^{2}(1+qC_\psi)}{\Lambda}
  \Bigl(\frac{\kappa_2C_U}{2\xi} + \frac43\Bigr)\Bigr\}.
\]
When $pq2^{J_n}\log(pq2^{J_n})/n \to 0$, $\lamGp = \lamG\{1 + o(1)\}$.

\begin{proposition}[Proposition 3 above the concentration ceiling]
\label{prop:slow-ceiling}
Let the assumptions of Proposition~3 hold, with the conditions $2^{J_n} \le 2n$
and $pq2^{J_n}\log(pq2^{J_n})/n \to 0$ replaced by $2^{J_n} \le n^{2}$, and
with $\lamGp$ of \eqref{eq:mu-blocks} in place of $\lamG$ in the estimator.
Then
\[
  \normn{\widehat{f}-f}^{2} = O_p\bigl\{\lamGp\nGw{\thetastar} +
  \mathcal{B}_{J_n}^{2} + p/n\bigr\},
\]
where $\lamGp\nGw{\thetastar} \le \bar\varrho\,(\mu^{\Gcal}_{J_n})^{1/2}
\lamGone\nGone{\thetastar}$,
with $(\lamGone)^{2} \le (C_{\pen} + 16\Lambda)\sigma^{2}\bn/n$ and $C_{\pen}$
as in \eqref{eq:order-lambda}, and:
\begin{enumerate}
\item[(iii)] (the pairs that Proposition~3(i) and (ii) exclude, all with
$\pi < 2$) (a) if $s > 1/2$ and $\seff \le 1/4$, the choice
$2^{J_n} \asymp (n/\log n)^{2/(1+4\seff)}$ gives
\[
  \normn{\widehat{f}-f}^{2} = O_p\bigl\{(\log n/n)^{4\seff/(1+4\seff)}\bigr\};
\]
(b) if $s < 1/2$ and $\seff \le s/2$, the choice
$2^{J_n} \asymp (n/\bn^{1/\pi})^{1/(1-s+2\seff)}$ gives
\[
  \normn{\widehat{f}-f}^{2}
  = O_p\bigl\{n^{-2\seff/(1-s+2\seff)}
  (\log n)^{(1/\pi)2\seff/(1-s+2\seff)}\bigr\};
\]
in both, $2^{J_n} \gtrsim n/\log n$, and the powers of $n$ join those of
Proposition~3(i) and (ii) at $\seff = 1/4$ and $\seff = s/2$;
\item[(iv)] (the boundary $s = 1/2$, with $\pi > 1$)
$\nGone{\thetastar} \le pqC_g\{(1-2^{-\seff})^{-1} + \bn^{-\vartheta_{\pi}}J_n\}$,
and: if $\seff \ge 1/4$, that is $\pi \ge 4/3$, the choice
$2^{J_n} \asymp (n/\bn)^{1/(4\seff)}$ gives
$\normn{\widehat{f}-f}^{2} = O_p\{n^{-1/2}(\log n)^{3/2-\vartheta_{\pi}}\}$,
which is $n^{-1/2}\log n$ when $\pi \ge 2$; if $\seff < 1/4$, that is
$1 < \pi < 4/3$, the choice
$2^{J_n} \asymp \{n/\bn^{1+1/\pi}\}^{2/(1+4\seff)}$ gives
$\normn{\widehat{f}-f}^{2} =
O_p[\{(\log n)^{1+1/\pi}/n\}^{4\seff/(1+4\seff)}]$.
\end{enumerate}
The proof uses no lower bound on $\Gramhat$, and no property of $\Gramhat$
other than the upper bound $\lmaxx(\Gramhat) \le \Lambda\mu^{\Gcal}_{J_n}$, which
holds with probability tending to one under the upper bounds of Assumptions~2
and~3 alone, with no condition on $J_n$.
\end{proposition}

Stopping $J_n$ at the ceiling instead, with $2^{J_n} = n/(K\log n)$ and $K$
fixed and large enough, keeps $\lamG$ admissible and gives
$O_p\{(\log n/n)^{2\seff}\}$ in the pairs of (iii); above the ceiling the
estimation term grows more slowly than the bias decreases, and the exponents
of (iii) are strictly larger than $2\seff$ away from $\seff = 1/4$ and
$\seff = s/2$ (Remark~\ref{rem:ceiling} in Section~S8, for the lasso). Compared
with Proposition~\ref{prop:lasso-slow}(iii) and (iv) for the lasso, the chunks
remove nothing in (iii)(a), where the coarse chunk dominates both norms, as in
Proposition~3(i); the factor $(\log n)^{(1-1/\pi)2\seff/(1-s+2\seff)}$ in
(iii)(b); and, in (iv), $(\log n)^{\vartheta_{\pi}}$ when $\seff \ge 1/4$, half
a logarithm when $\pi \ge 2$, and $(\log n)^{(1-1/\pi)4\seff/(1+4\seff)}$ when
$\seff < 1/4$. This is the mechanism of Proposition~3(ii): the logarithm comes
from the size of the chunks, and the count of Lemma~\ref{lem:count} returns it
to the extent that the signal is spread over them.

\subsection*{Proof of Proposition~\ref{prop:slow-ceiling}}

\emph{The two events above the ceiling.} Let
$\Acal^{++} = \{\smax^{2} \le 2B_X^{2}C_U\mu^{\Gcal}_{J}\}\cap
\{\lmaxx(\Gramhat) \le \Lambda\mu^{\Gcal}_{J}\}$. For the first, the Bernstein
bound of Lemma~\ref{lem:columns} holds for every $J$ and every $t > 0$. With
$x = 2\log d$ and $t = (2xR'_JB_X^{2}C_U/n)^{1/2} + 2xR'_J/(3n)$,
$nt^{2} \ge 2xR'_JB_X^{2}C_U + (2/3)xR'_Jt$, so that the exponent is at most
$-x$; and $(2ab)^{1/2} \le a + b/2$ with $a = B_X^{2}C_U$ and $b = xR'_J/n$
gives $t \le B_X^{2}C_U + (7/6)xR'_J/n$. Since
$2B_X^{2}C_U + (7/3)R'_J\log d/n = 2B_X^{2}C_U(1 + c_{\psi}2^{J}\log d/n)$,
\begin{equation}
  \label{eq:smax-ceiling}
  \Prob\bigl\{\smax^{2} > 2B_X^{2}C_U(1 + c_{\psi}2^{J}\log d/n)\bigr\}
  \;\le\; \frac{1}{d}
  \qquad\text{for every } J,
\end{equation}
and $c_{\psi}2^{J}\log d \le c_{\mu}2^{J}\log(2D)$. For the second, the matrix
Bernstein bound in the proof of Proposition~2(ii),
$\Prob(\opnorm{\Gramhat-\Gram} \ge t) \le
2D\exp\{-nt^{2}/(2R_J(\kappa_2C_U + 2t/3))\}$, holds for every $J$ and every
$t > 0$. With $x = 2\log(2D)$ and
$t = (2xR_J\kappa_2C_U/n)^{1/2} + (4/3)xR_J/n$,
$nt^{2} \ge 2xR_J\kappa_2C_U + (4/3)xR_Jt$, and the probability is at most
$2De^{-x} = 1/(2D)$; and $(2ab)^{1/2} \le \xi + ab/(2\xi)$ with
$a = \kappa_2C_U$ and $b = xR_J/n$ gives
$t \le \xi + \{\kappa_2C_U/(2\xi) + 4/3\}xR_J/n$. By Weyl's inequality and
$\lmaxx(\Gram) \le \kappa_2C_U$, off that event
$\lmaxx(\Gramhat) < \Lambda + \{\kappa_2C_U/(2\xi) + 4/3\}2\log(2D)R_J/n \le
\Lambda\mu^{\Gcal}_{J}$, because
$R_J = pB_X^{2}(1 + qC_\psi2^{J}) \le pB_X^{2}(1+qC_\psi)2^{J}$. Hence
$\Prob(\Acal^{++}) \ge 1 - 1/d - 1/(2D) \to 1$, with no condition on $J_n$.
Both bounds use only $|X_{\cv}| \le B_X$, $p_{\bU} \le C_U$ and
$\lmaxx(\Gram) \le \kappa_2C_U$.

\emph{The general bound.} On $\Acal^{++}$ the two bounds in the proof of
Lemma~\ref{lem:calib-blocks}(ii), $\mathrm{tr}(\Psitil_{G}) \le |G|\smax^{2}$
and $\opnorm{\Psitil_{G}} \le \lmaxx(\Gramhat)$, give
$\lamz{G} \le (\mu^{\Gcal}_{J})^{1/2}\lamzbar{G}$ for every $G$, so that
$\lamGp \ge \pen_{w}$ and $\Ecal \subseteq \Tw$, and the proof of
Proposition~3 goes through with $\lamGp$ in place of $\lamG$ and $\Acal^{++}$
in place of $\Acal^{+}$; the chain gives
$\lamGp\nGw{\thetastar} \le \bar\varrho(\mu^{\Gcal}_{J_n})^{1/2}
\lamGone\nGone{\thetastar}$. With $2^{J_n} \le n^{2}$,
$|\Gcal_{n}| \le pq(1 + n^{2}/\bn) \le pqn^{2}$ and
$\log(|\Gcal_{n}|/\alpha) \le \log(2pq/\alpha) + 2\bn$, and the computation in
the proof of Lemma~\ref{lem:regime} gives
$(\lamGone)^{2} \le (16\sigma^{2}/n)\{(2B_X^{2}C_U + 2\Lambda)\bn +
\Lambda\log(2pq/\alpha)\} \le (C_{\pen} + 16\Lambda)\sigma^{2}\bn/n$. Since
$\log(2D_n) \asymp \log n \asymp \bn$,
\begin{equation}
  \label{eq:lam-plus-blocks}
  (\mu^{\Gcal}_{J_n})^{1/2}\lamGone \;\asymp\;
  \Bigl(\frac{\log n}{n}\Bigr)^{1/2}
  \Bigl(1 + \frac{2^{J_n}\log n}{n}\Bigr)^{1/2},
\end{equation}
which is of order $2^{J_n/2}\log n/n$ when $2^{J_n} \gtrsim n/\log n$.

\emph{Part (iii)(a).} With $s > 1/2$, $\nGone{\thetastar} = O(1)$ as in the
proof of Proposition~3(i). With
$2^{J_n} \asymp (n/\log n)^{2/(1+4\seff)} \gtrsim n/\log n$, because
$\seff \le 1/4$, \eqref{eq:lam-plus-blocks} gives the estimation term
$(\log n/n)(n/\log n)^{1/(1+4\seff)} = (\log n/n)^{4\seff/(1+4\seff)}$, the
order of the bias $2^{-2J_n\seff}$; $p/n$ is of smaller order.

\emph{Part (iii)(b).} With $s < 1/2$ and $\pi < 2$,
$\vartheta_{\pi} = 1 - 1/\pi$ and
$\nGone{\thetastar} = O\{1 + \bn^{-\vartheta_{\pi}}2^{J_n(1/2-s)}\}$ as in the
proof of Proposition~3(ii). With
$2^{J_n} \asymp (n/\bn^{1/\pi})^{1/(1-s+2\seff)}$, which exceeds $n/\bn$ for
$n$ large, because the exponent is at least $1$ and $1/\pi \le 1$, and which is
at most $n^{2}$, because $1-s+2\seff > 1/2$, \eqref{eq:lam-plus-blocks} gives
the estimation term
\[
  2^{J_n/2}\frac{\log n}{n}\bigl\{1 + \bn^{-(1-1/\pi)}2^{J_n(1/2-s)}\bigr\}
  \;\asymp\; 2^{J_n/2}\frac{\log n}{n} +
  2^{J_n(1-s)}\frac{(\log n)^{1/\pi}}{n},
\]
whose second term dominates and, at the choice of the statement, equals
$(\bn^{1/\pi}/n)^{1-(1-s)/(1-s+2\seff)} =
(\bn^{1/\pi}/n)^{2\seff/(1-s+2\seff)}$, the order of the bias
$2^{-2J_n\seff}$.

\emph{Part (iv).} With $s = 1/2$, the sum in \eqref{eq:count} is at most
$J_n$, which gives the bound on $\nGone{\thetastar}$ of the statement, and
$J_n \asymp \log n \asymp \bn$ in both choices, so that
$\nGone{\thetastar} \asymp (\log n)^{1-\vartheta_{\pi}}$, since
$\vartheta_{\pi} \le 1/2$. If $\seff \ge 1/4$,
$2^{J_n} \asymp (n/\bn)^{1/(4\seff)} \lesssim n/\log n$,
$\mu^{\Gcal}_{J_n}$ is bounded, the estimation term is
$(\log n/n)^{1/2}(\log n)^{1-\vartheta_{\pi}}$, and the bias, of order
$(\bn/n)^{1/2}$, is of smaller order. If $\seff < 1/4$, then $\pi < 4/3$ and
$\vartheta_{\pi} = 1 - 1/\pi$; with
$2^{J_n} \asymp (n/\bn^{1+1/\pi})^{2/(1+4\seff)}$, which exceeds $n/\log n$,
the estimation term is $2^{J_n/2}(\log n)^{1+1/\pi}/n \asymp
\{(\log n)^{1+1/\pi}/n\}^{4\seff/(1+4\seff)}$, the order of the bias. \qed

% derivations/08-blocos.tex, prova do Teorema 4 (numeracao global).
\subsection*{Proof of Theorem 3}

By the definition of $J_n$, $n^{1/(2\seff+1)} \le 2^{J_n} <
2n^{1/(2\seff+1)} \le 2n$, so the setting of this section holds, and
$pq2^{J_n}\log(pq2^{J_n})/n \asymp n^{1/(2\seff+1)-1}\log n =
n^{-2\seff/(2\seff+1)}\log n \to 0$, which is the requirement of
Proposition~2(ii); by Lemma~\ref{lem:regime}, $\Prob(\Acal) \to 1$. Fix
$\epsilon > 0$, take $\alpha = \alpha' = \alpha'' = \epsilon/4$, and intersect
$\Acal$, the event $\Ecal$ of Lemma~\ref{lem:calib-blocks}(i),
$\{\normn{\bb}^{2} \le \mathcal{B}_{J_n}^{2}/\alpha'\}$ (Markov's inequality,
since $\E\normn{\bb}^{2} = \norm{f-f_{J_n}}_{L_2(P)}^{2} \le
\mathcal{B}_{J_n}^{2}$ by Lemma~2) and
$\{\normn{\PA\bveps}^{2} \le \sigma^{2}p/(n\alpha'')\}$ (Markov's inequality and
Theorem~2(iii)); for $n$ large the intersection has probability at least
$1-\epsilon$. On it $\Tw$ holds, and \eqref{eq:oracle-min-blocks} with
$\bar\bbeta = \betastar$ gives
\[
  \normn{\widehat{f}-f}^{2} \;\le\;
  \frac{192(\lamG)^{2}\Wcal(\Gzero(\thetastar))}{\gtil}
  + 60\normn{\bb}^{2} + 3\normn{\PA\bveps}^{2}.
\]
By Lemma~\ref{lem:calib-blocks}(iii), Lemma~\ref{lem:regime} and
Proposition~\ref{prop:balanced}(ii), with $\gtil > \gamma/2$,
\[
  \frac{192(\lamG)^{2}\Wcal(\Gcal_{n})}{\gtil}
  \;\le\; \frac{384\bar\varrho^{2}(\lamGone)^{2}|\Gcal_{n}|}{\gamma}
  \;\le\; \frac{384\bar\varrho^{2}C_{\pen}\sigma^{2}pq(\bn + 2^{J_n})}{\gamma n},
\]
which is of order $n^{-2\seff/(2\seff+1)}$, because
$2^{J_n} \asymp n^{1/(2\seff+1)}$ dominates $\bn \asymp \log n$. The bias term
is $60\mathcal{B}_{J_n}^{2}/\alpha' \asymp 2^{-2J_n\seff} \le
n^{-2\seff/(2\seff+1)}$, by the lower bound on $2^{J_n}$, and the last term is
$O(1/n)$, of strictly smaller order. Hence
$\Prob\{\normn{\widehat{f}-f}^{2} > M_\epsilon n^{-2\seff/(2\seff+1)}\} <
\epsilon$ for a constant $M_\epsilon$. \qed

% derivations/08-blocos.tex, prova do Corolario 10 (numeracao global).
\subsection*{Proof of Corollary 1}

By the orthonormality of the wavelets of level below $J_n$ in $L_2[0,1]$ and
Lemma~2,
$\sum_{\cv,\md}\lVert\widehat{g}_{\cv\md}-\gcomp{\cv}{\md}\rVert^{2}_{L_2[0,1]}
\le 2\norm{\thetahat-\thetastar}_{2}^{2} +
2pqC_g^{2}(1-2^{-2\seff})^{-1}2^{-2J_n\seff}$. On the event of the proof of
Theorem~3, Theorem~2(ii) with $\bar\bbeta = \betastar$ gives
$\norm{\thetahat-\thetastar}_{2}^{2} \le
64(\lamG)^{2}\Wcal(\Gcal_{n})/\gtil^{2} + 16\normn{\bb}^{2}/\gtil$, with
$\gtil > \gamma/2$: the two terms have the order of those of Theorem~3, with
$\gamma^{-2}$ in place of $\gamma^{-1}$. For the levels,
Lemma~\ref{lem:profile} with $\bar\bbeta = \betastar$ gives
$\Amat(\chat-\bc) = \PA(\bb+\bveps) - \PA\Bmat v$, so that
$\normn{\Amat(\chat-\bc)} \le \normn{\bb} + \normn{\PA\bveps} +
\Lambda^{1/2}\norm{v}_{2}$, and $\lmin(\Amat\T\Amat/n) > \gamma/2$ gives
$\norm{\chat-\bc}_{2}^{2} \le 2\normn{\Amat(\chat-\bc)}^{2}/\gamma$. Finally the
components of a functional coefficient depend on distinct coordinates of $u$
and the basis is orthonormal, so that
$\lVert\widehat{\fcoef{\cv}}-\fcoef{\cv}\rVert_{L_2([0,1]^q)} \le
|\widehat{c}_{\cv}-\lvl{\cv}| + \sum_{\md}\lVert\widehat{g}_{\cv\md}-
\gcomp{\cv}{\md}\rVert_{L_2[0,1]}$, with $q$ fixed. \qed

% derivations/08-blocos.tex, prova do Corolario 11 (numeracao global); D47(a).
\subsection*{Proof of Theorem 1}

\emph{The window.} Since $\seff \le s$,
$s/\{(2s+1)\seff\} \ge 1/(2s+1)$, and the window is nonempty if and only if
$s/\{(2s+1)\seff\} < 1$, that is $\seff > s/(2s+1)$. With $c < 1$,
$2^{J_n} \le 2n^{c} \le 2n$, so the setting of this section holds; the
requirement of Proposition~2(ii) is $pq2^{J_n}\log(pq2^{J_n})/n \asymp
n^{c-1}\log n \to 0$, and $\Prob(\Acal) \to 1$ by Lemma~\ref{lem:regime}.

\emph{The rate.} Fix $\epsilon > 0$ and the event of the proof of Theorem~3.
On it, Corollary~\ref{cor:ideal} with $\Lambda' = \Lambda$ gives
$\normn{\widehat{f}-f}^{2} \le 120\Lambda\Rid(\thetastar;\eta) +
120\normn{\bb}^{2} + 3\normn{\PA\bveps}^{2}$, and, by
Lemma~\ref{lem:calib-blocks}(iii), \eqref{eq:order-lambda} and
$\gtil > \gamma/2$,
\[
  \eta \;=\; \frac{1.6\max_{G}(\lamG w_{G})^{2}}{\Lambda\gtil}
  \;\le\; \frac{3.2\bar\varrho^{2}(\lamGone)^{2}}{\Lambda\gamma}
  \;\le\; \bar\eta_{n} := \kappa\,\frac{\bn}{n},
  \qquad
  \kappa = \frac{3.2\bar\varrho^{2}C_{\pen}\sigma^{2}}{\Lambda\gamma}.
\]
Since $\Rid$ is nondecreasing in $\eta$, Lemma~3 with $b = \bn$ gives
\[
  \Rid(\thetastar;\eta) \;\le\; \Rid(\thetastar;\bar\eta_{n})
  \;\le\; pq\Bigl\{\AG C_g^{\tau}\kappa^{2s/(2s+1)}\,n^{-2s/(2s+1)}
  \bn^{(2/\pi-1)_{+}/(2s+1)} + \kappa\frac{\bn}{n}\Bigr\},
\]
with $\bn \asymp \log n$; the second term is of smaller order because
$2s/(2s+1) < 1$. The bias is $O(2^{-2J_n\seff}) = O(n^{-2c\seff})$, and
$2c\seff \ge 2s/(2s+1)$ by the window; the term $\sigma^{2}p/n$ is of smaller
order. For the components, the second bound of Corollary~\ref{cor:ideal} and
Lemma~2 give the same rate, as in the proof of Corollary~1.

% E5g (pergunta 51(c), D57): a uniformidade do O_p, lida na prova do
% Corolario 16 de derivations/09-cota-inferior.tex; sem ela, "minimax optimal"
% no corpo fica no sentido ponto a ponto.
\emph{Uniformity over the class.} The events $\Acal$ and $\Ecal$ and the
event for $\normn{\PA\bveps}^{2}$ do not depend on $f$; the event for
$\normn{\bb}^{2}$ does, but its probability is at least $1-\alpha'$ for every
$f$ whose components satisfy Assumption~5, because the bound
$\mathcal{B}_{J_n}^{2}$ on $\E\normn{\bb}^{2}$ of Lemma~2 depends on $f$ only
through $C_g$. On the event, the bound depends on $f$ only through $C_g$, by
Lemmas~2 and~3, and not on the levels $\bc$, since $\bb = f - f_{J_n}$ does
not involve them. Hence
$\sup_{f}\Prob_{f}\{\normn{\widehat{f}-f}^{2} > M_\epsilon r_n\} < \epsilon$
for $n$ large, with $\Prob_{f}$ the law of the sample when the regression
function is $f$ and the supremum over the functions of Assumption~5 with any
levels, and the same holds for the components: the bound of Theorem~1 is
uniform over the class, which is the form that Theorem~\ref{thm:lower} below
is compared with.

\emph{Two readings.} The constant grows with the weights at most as
$\bar\varrho^{\,4s/(2s+1)}$ in the leading term, through
$\kappa^{2s/(2s+1)}$. The gain over the rate $(\log n/n)^{2s/(2s+1)}$ of the
lasso (Theorem~\ref{thm:lasso-main}) is the factor $(\log n)^{2\seff/(2s+1)}$,
because $2s/(2s+1) - (2/\pi-1)_{+}/(2s+1) = 2\seff/(2s+1)$: when $\pi < 2$,
$2s - 2/\pi + 1 = 2(s - 1/\pi + 1/2) = 2\seff$, and when $\pi \ge 2$,
$\seff = s$. The logarithm of the union bound over the columns is absorbed by
the size of the chunks, which is why they must have at least of the order of
$\log|\Gcal_{n}|$ coefficients. \qed

% derivations/08-blocos.tex, Observacao rem:inferior.
% E5g (D57): removida inteira e substituida pelo Theorem S6.1 abaixo; fica no
% lugar em cinza e sem numero (D51(b)), sem o rotulo rem:lower, que nada cita.


% E5g (E1.9; D57): derivations/09-cota-inferior.tex, Teorema 5 e Lema 17
% (numeracao global), nas tres perdas, com as provas; dentro do grupo azul de
% k = 4. Os simbolos locais do 09 (o hipercubo C, Omega, zeta, d_H, a funcao
% de limiar) ficam escritos no texto; |C| no lugar do M do 09, que colide com
% M_epsilon e com o M da S8; L_n e L_P escritos por extenso, porque L_n e o
% logaritmo da S8.
\subsection*{A minimax lower bound}

The rates of this section are upper bounds. For a lower bound, fix any law
$P$ of $(\bX,\bU)$ that satisfies Assumptions~2 and~3, with $\bX$ allowed to
depend on $\bU$, and let the errors be Gaussian,
\begin{equation}
  \label{eq:gauss}
  \varepsilon_i \sim N(0,\sigma^{2}),\ \text{independent of } (\bX_i,\bU_i),
\end{equation}
which satisfy Assumption~1. Write $\Prob_f$ and $\E_f$ for the law and the
expectation of the sample when the regression function is $f$, and
$\mathcal{F} = \mathcal{F}(s,\pi,C_g)$ for the set of the functions
$f(x,u) = \sum_{\cv}x_{\cv}\{\lvl{\cv} + \sum_{\md}\gcomp{\cv}{\md}(u_{\md})\}$
with $\bc \in \R^{p}$ arbitrary and every component satisfying the centering
and the bound of Assumption~5; by Proposition~1, $f$ determines its levels and
its components. An estimator is any measurable function of the sample: for
prediction, $\widehat{f}$ with values in $L_2(P)$, or only its values
$\widehat{f}(\bX_i,\bU_i)$, and for the components,
$\widehat{g} = (\widehat{g}_{\cv\md})$ with $\widehat{g}_{\cv\md} \in
L_2[0,1]$. For a nonempty set $\Kcal$ of blocks $(\cv,\md)$, the loss for the
components in $\Kcal$ is
\[
  L_{\Kcal}(\widehat{g},f) \;=\; \sum_{(\cv,\md)\in\Kcal}
  \lVert\widehat{g}_{\cv\md} - \gcomp{\cv}{\md}\rVert_{L_2[0,1]}^{2} .
\]
With $\gamma = \kappa_1c_U$ as in Section~S1 and
$\cA = (1-2^{-1/2})/2 \approx 0.146$, put, for
$n \ge n_{0} = \lceil 2\sigma^{2}/(C_g^{2}B_X^{2}C_U)\rceil$,
\begin{equation}
  \label{eq:lower-choices}
  \begin{gathered}
  \zeta_n^{2} = \frac{2\sigma^{2}}{nB_X^{2}C_U},
  \qquad
  \lev_n = \max\Bigl\{\lev \ge 0 : 2^{\lev(2s+1)} \le \frac{C_g^{2}}{\zeta_n^{2}}\Bigr\},
  \\
  \rlow^{2} = C_g^{2/(2s+1)}\Bigl(\frac{\sigma^{2}}{nB_X^{2}C_U}\Bigr)^{2s/(2s+1)} ;
  \end{gathered}
\end{equation}
$\lev_n$ exists because $C_g^{2}/\zeta_n^{2} \ge 1$ when $n \ge n_{0}$, and
$\rlow^{2}$, of order $n^{-2s/(2s+1)}$, is the rate of the lower bound.

\begin{theorem}[Minimax lower bound]
\label{thm:lower}
Let $P$ be any law of $(\bX,\bU)$ that satisfies Assumptions~2 and~3, with
$q \ge 1$ and $\bX$ possibly dependent on $\bU$, let the basis satisfy
Assumption~4, and let the errors be Gaussian, as in \eqref{eq:gauss}. For
every $s > 0$, every $1 \le \pi \le \infty$, every nonempty set $\Kcal$ of
blocks and every $n \ge n_{0}$, with $\rlow$ as in \eqref{eq:lower-choices}:
\begin{enumerate}
\item[(i)] (components)
\begin{gather*}
  \inf_{\widehat{g}}\sup_{f\in\mathcal{F}}\E_{f}L_{\Kcal}(\widehat{g},f)
  \;\ge\; \frac{\cA}{8}\,|\Kcal|\,\rlow^{2},
  \\
  \inf_{\widehat{g}}\sup_{f\in\mathcal{F}}\Prob_{f}\Bigl\{L_{\Kcal}(\widehat{g},f)
  \ge \frac{\cA}{16}\,|\Kcal|\,\rlow^{2}\Bigr\} \;\ge\; \frac{\cA}{2} ;
\end{gather*}
\item[(ii)] (prediction, population norm) the same holds for
$\norm{\widehat{f}-f}_{L_2(P)}^{2}$, with $\gamma|\Kcal|\rlow^{2}$ in place of
$|\Kcal|\rlow^{2}$;
\item[(iii)] (prediction, empirical norm) if moreover $n \ge n_{1}$, the
smallest integer from which on
$2D\exp[-n\gamma^{2}/\{8R_{J}(\kappa_2C_U + \gamma/3)\}] \le \cA/4$, with
$J = \lev_n + 1$ and $D = p + pq\NJ$ and $R_J$ as in Proposition~2(ii), then
\begin{gather*}
  \inf_{\widehat{f}}\sup_{f\in\mathcal{F}}\E_{f}\normn{\widehat{f}-f}^{2}
  \;\ge\; \frac{\cA}{32}\,\gamma|\Kcal|\,\rlow^{2},
  \\
  \inf_{\widehat{f}}\sup_{f\in\mathcal{F}}\Prob_{f}\Bigl\{\normn{\widehat{f}-f}^{2}
  \ge \frac{\cA}{32}\,\gamma|\Kcal|\,\rlow^{2}\Bigr\} \;\ge\; \frac{\cA}{4} .
\end{gather*}
\end{enumerate}
The constants depend neither on $\pi$, $r$, $p$ or $q$, nor on $P$ beyond
$(B_X,C_U)$ in $\rlow$ and $\gamma = \kappa_1c_U$; $n_{0}$ depends only on
$(\sigma,C_g,B_X,C_U)$, and $n_{1}$, which is finite, also on
$(p,q,s,\kappa_1,\kappa_2,c_U)$ and the wavelet family. In particular, with
$\Kcal$ a single block, every component and the prediction have minimax risk
of order at least $n^{-2s/(2s+1)}$, and with $\Kcal$ all the blocks, so that
$|\Kcal| = pq$, the bound is proportional to $pq$.
\end{theorem}

The bound is not asymptotic, and it holds for every pair $(s,\pi)$, in
particular for every pair that Assumption~5 allows. It has the form of the
lower bound of \citet[Theorem~5]{Donoho-Johnstone-1998} for Besov bodies in
the Gaussian sequence model, with radius $C_g$ and noise level
$\sigma(nB_X^{2}C_U)^{-1/2}$ per coordinate. Since Gaussian errors satisfy
Assumption~1, it holds a fortiori over the models of Assumptions~1 to~5, and
the form in probability is the one to compare with the upper bounds, which are
stated in $O_p$ and are uniform over the class (proof of Theorem~1 above).
Hence, when $\pi \ge 2$, $n^{-2s/(2s+1)}$ is the minimax rate in the
prediction norm and for the components, for every $q$ and with covariates
that depend on the modulators, and the estimator of Theorem~1 attains it
uniformly over the class; when $\pi < 2$ and $\seff > s/(2s+1)$ it attains it
up to the factor $(\log n)^{(2/\pi-1)/(2s+1)}$, which the lower bound does not
carry
% D58 (pergunta 52(f)): a condicao de Cai (1999), que saiu com a Remark S6.1.
(the same factor as in the rate of block thresholding in the sequence model,
\citealp[Theorem~4]{Cai-1999}, which holds for $s \ge 1/\pi$, a condition
strictly stronger than $\seff > s/(2s+1)$), and the coordinatewise lasso, at the rate $(\log n/n)^{2s/(2s+1)}$ of
Theorem~\ref{thm:lasso-main}, stays above it by the factor
$(\log n)^{2s/(2s+1)}$ for every $\pi$. When $q = 1$ and $\bX$ is independent
of $U_1$, part (i) with $\Kcal$ all the blocks is, up to the constants, the
nonparametric term of the lower bound of
\citet[Theorem~1, eqs.~3.6 and~3.7]{Klopp-Pensky-2015} with a fixed number of
covariates.

The proof applies Assouad's lemma to a hypercube of wavelets at a single
level. Fix $\Kcal$, a level $\lev \ge 0$ and a height $\zeta > 0$; let
\[
  \mathcal{C} = \bigl\{(\cv,\md,\tsl) : (\cv,\md)\in\Kcal,\ 0 \le \tsl < 2^{\lev}\bigr\},
  \qquad |\mathcal{C}| = |\Kcal|\,2^{\lev},
\]
and $\Omega = \{0,1\}^{\mathcal{C}}$. For $\omega \in \Omega$ the components
are $\gcomp{\cv}{\md}^{\omega} =
\zeta\sum_{\tsl}\omega_{\cv\md,\tsl}\wav{\lev}{\tsl}$ for
$(\cv,\md) \in \Kcal$ and zero otherwise, the levels are zero, and
\[
  f_{\omega}(x,u) \;=\; \zeta\,z_{\mathcal{C}}(x,u)\T\omega,
  \qquad
  z_{\mathcal{C}}(x,u) = \bigl(x_{\cv}\wav{\lev}{\tsl}(u_{\md})\bigr)_{(\cv,\md,\tsl)\in\mathcal{C}} ,
\]
whose entries are columns of the design at resolution $J = \lev + 1$. Write
$\Gram_{\mathcal{C}\mathcal{C}}$ and $\Gramhat_{\mathcal{C}\mathcal{C}}$ for
the principal submatrices of $\Gram$ and $\Gramhat$ at that resolution on the
columns in $\mathcal{C}$, $\bZ_{\mathcal{C}}$ for those columns of $\bZ$,
$d_{H}(\omega,\omega') = \sum_{(\cv,\md,\tsl)\in\mathcal{C}}
\mathbf{1}\{\omega_{\cv\md,\tsl} \ne \omega'_{\cv\md,\tsl}\}$ for the Hamming
distance, and $\Prob_{\omega}$ and $\E_{\omega}$ for $\Prob_{f_{\omega}}$ and
$\E_{f_{\omega}}$. The form of Assouad's lemma that we use is
\citet[Theorem~2.12(iv)]{Tsybakov-2009} with the Kullback--Leibler divergence
bounded by one: if $\operatorname{KL}(\Prob_{\omega},\Prob_{\omega'}) \le 1$
whenever $d_{H}(\omega,\omega') = 1$, then
\begin{equation}
  \label{eq:assouad}
  \inf_{\widehat\omega}\max_{\omega\in\Omega}\E_{\omega}d_{H}(\widehat\omega,\omega)
  \;\ge\; \frac{|\mathcal{C}|}{2}\max\Bigl(\frac{e^{-1}}{2},\,1-2^{-1/2}\Bigr)
  \;=\; \cA|\mathcal{C}| ,
\end{equation}
the infimum being over all the estimators with values in $\Omega$.

\begin{lemma}[A hypercube of wavelets]
\label{lem:hypercube}
In the setting of Theorem~\ref{thm:lower}, for every nonempty $\Kcal$, every
$\lev \ge 0$ and every $\zeta > 0$:
\begin{enumerate}
\item[(i)] $f_{\omega} \in \mathcal{F}$ for every $\omega \in \Omega$ if and
only if $\zeta2^{\lev(s+1/2)} \le C_g$, for every $s > 0$ and every
$1 \le \pi \le \infty$;
\item[(ii)] for every pair,
$\operatorname{KL}(\Prob_{\omega},\Prob_{\omega'}) =
\{n\zeta^{2}/(2\sigma^{2})\}(\omega-\omega')\T\Gram_{\mathcal{C}\mathcal{C}}(\omega-\omega')$,
and, if $\omega$ and $\omega'$ differ only in the coordinate $(\cv,\md,\tsl)$,
\[
  \operatorname{KL}(\Prob_{\omega},\Prob_{\omega'})
  \;=\; \frac{n\zeta^{2}}{2\sigma^{2}}\,\E\bigl\{X_{\cv}^{2}\wav{\lev}{\tsl}(U_{\md})^{2}\bigr\}
  \;\le\; \frac{n\zeta^{2}}{2\sigma^{2}}\,\min(B_X^{2},\kappa_2)\,C_U ;
\]
\item[(iii)] $\lmin(\Gram_{\mathcal{C}\mathcal{C}}) \ge \gamma$, so that
\begin{gather*}
  \norm{f_{\omega}-f_{\omega'}}_{L_2(P)}^{2} \ge \gamma\zeta^{2}d_{H}(\omega,\omega'),
  \\
  \sum_{(\cv,\md)\in\Kcal}\lVert\gcomp{\cv}{\md}^{\omega} -
  \gcomp{\cv}{\md}^{\omega'}\rVert_{L_2[0,1]}^{2} = \zeta^{2}d_{H}(\omega,\omega') ;
\end{gather*}
with a single block $(\cv,\md)$, the first statement uses only
$\E(X_{\cv}^{2}\mid\bU) \ge \kappa_1$ and the lower bound $c_U$ on the
density of $U_{\md}$;
\item[(iv)] for each of the three losses and every estimator there is a
measurable function $\widehat\omega$ of the sample with values in $\Omega$
such that, for every $\omega \in \Omega$,
\begin{enumerate}
\item[(a)] $L_{\Kcal}(\widehat{g},f_{\omega}) \ge (\zeta^{2}/4)\,d_{H}(\widehat\omega,\omega)$,
\item[(b)] $\norm{\widehat{f}-f_{\omega}}_{L_2(P)}^{2} \ge (\gamma\zeta^{2}/4)\,d_{H}(\widehat\omega,\omega)$,
\item[(c)] $\normn{\widehat{f}-f_{\omega}}^{2} \ge (\gamma\zeta^{2}/8)\,
d_{H}(\widehat\omega,\omega)\,\mathbf{1}_{\mathcal{E}_{\mathcal{C}}}$, where
$\mathcal{E}_{\mathcal{C}} = \{\lmin(\Gramhat_{\mathcal{C}\mathcal{C}}) \ge \gamma/2\}$
depends only on the $(\bX_i,\bU_i)$ and has, with $J = \lev+1$ and $D$ and
$R_J$ as in Theorem~\ref{thm:lower}(iii),
\[
  \Prob(\mathcal{E}_{\mathcal{C}}^{c}) \;\le\;
  2D\exp\Bigl\{-\frac{n\gamma^{2}}{8R_{J}(\kappa_2C_U + \gamma/3)}\Bigr\} .
\]
\end{enumerate}
\end{enumerate}
\end{lemma}

The Kullback--Leibler divergence uses only upper bounds on the design, and
through the conditional second moment $\E(X_{\cv}^{2}\mid\bU)$, which is why
the covariates may depend on the modulators at no cost. The separation of the
components uses no property of the design; that of the prediction uses the
smallest eigenvalue, and it is there that the cross terms between modulators
and between covariates enter, through Proposition~2(i).

\subsection*{Proof of Lemma~\ref{lem:hypercube}}

\emph{Part (i).} By the orthonormality of $\{\mathbf{1}\}\cup\{\wav{\lev}{\tsl}\}$
in $L_2[0,1]$ (Section~S2), the coefficients of $\gcomp{\cv}{\md}^{\omega}$
are $\theta_{\cv\md,\lev'\tsl} = \zeta\omega_{\cv\md,\tsl}\mathbf{1}\{\lev' =
\lev\}$, and $\int_0^1\gcomp{\cv}{\md}^{\omega} = 0$ because
$\int_0^1\wav{\lev}{\tsl} = 0$. Only the level $\lev$ is nonzero, so the bound
of Assumption~5 reduces to that level, whatever the index $r$. For
$\pi < \infty$, $\norm{\theta_{\cv\md,\lev\cdot}}_{\pi} =
\zeta(\sum_{\tsl}\omega_{\cv\md,\tsl})^{1/\pi} \le \zeta2^{\lev/\pi}$, with
equality at $\omega \equiv 1$; for $\pi = \infty$ it is $\zeta$ if
$\omega_{\cv\md,\cdot} \ne 0$. In both cases the maximum over $\Omega$ of
$2^{\lev(s+1/2-1/\pi)}\norm{\theta_{\cv\md,\lev\cdot}}_{\pi}$ is
$\zeta2^{\lev(s+1/2)}$. The components outside $\Kcal$ and the levels are
zero.

\emph{Part (ii).} Under $\Prob_{\omega}$ the pairs $(\bX_i,\bU_i)$ have the
law $P$, the same for every $\omega$, and, given them, the $Y_i$ are
independent with $Y_i \sim N\{f_{\omega}(\bX_i,\bU_i),\sigma^{2}\}$. The
divergence between product laws is the sum of the divergences, and that
between two laws with the same marginal is the mean of the divergence between
the conditionals, so that
$\operatorname{KL}(\Prob_{\omega},\Prob_{\omega'}) =
n\E\{(f_{\omega}-f_{\omega'})^{2}(\bX,\bU)\}/(2\sigma^{2})$, and
$f_{\omega}-f_{\omega'} = \zeta z_{\mathcal{C}}\T(\omega-\omega')$ gives the
quadratic form. For a neighbor, $\E\{(f_{\omega}-f_{\omega'})^{2}\} =
\zeta^{2}\E\{\E(X_{\cv}^{2}\mid\bU)\wav{\lev}{\tsl}(U_{\md})^{2}\}$, with
$\E(X_{\cv}^{2}\mid\bU) \le B_X^{2}$ and $\E(X_{\cv}^{2}\mid\bU) =
e_{\cv}\T\E(\bX\bX\T\mid\bU)e_{\cv} \le \kappa_2$ by Assumption~2, and
$\E\{\wav{\lev}{\tsl}(U_{\md})^{2}\} \le C_U\int_0^1\wav{\lev}{\tsl}^{2} = C_U$,
because the density of $U_{\md}$ inherits the bound $C_U$ of Assumption~3.

\emph{Part (iii).} For $v \in \R^{\mathcal{C}}$ and $\tilde v$ its completion
with zeros, $v\T\Gram_{\mathcal{C}\mathcal{C}}v = \tilde v\T\Gram\tilde v \ge
\lmin(\Gram)\norm{v}_{2}^{2} \ge \gamma\norm{v}_{2}^{2}$, by
Proposition~2(i). Since $\omega-\omega' \in \{-1,0,1\}^{\mathcal{C}}$,
$\norm{\omega-\omega'}_{2}^{2} = d_{H}(\omega,\omega')$, and
$\norm{f_{\omega}-f_{\omega'}}_{L_2(P)}^{2} =
\zeta^{2}(\omega-\omega')\T\Gram_{\mathcal{C}\mathcal{C}}(\omega-\omega') \ge
\gamma\zeta^{2}d_{H}(\omega,\omega')$. The identity for the components is
Parseval's. With a single block,
\begin{align*}
  v\T\Gram_{\mathcal{C}\mathcal{C}}v &=
  \E\Bigl[\E(X_{\cv}^{2}\mid\bU)\Bigl\{\sum_{\tsl}v_{\tsl}\wav{\lev}{\tsl}(U_{\md})\Bigr\}^{2}\Bigr]
  \\
  &\ge \kappa_1c_U\int_0^1\Bigl(\sum_{\tsl}v_{\tsl}\wav{\lev}{\tsl}\Bigr)^{2} =
  \gamma\norm{v}_{2}^{2} .
\end{align*}

\emph{Part (iv).} Let $\varpi(y) = \mathbf{1}\{y \ge \zeta/2\}$ for real $y$.
For $\omega_{\cv\md,\tsl} \in \{0,1\}$,
\begin{equation}
  \label{eq:hypercube-threshold}
  |y - \zeta\omega_{\cv\md,\tsl}| \;\ge\; \frac{\zeta}{2}\,
  |\varpi(y) - \omega_{\cv\md,\tsl}| ,
\end{equation}
because if $\varpi(y) \ne \omega_{\cv\md,\tsl}$ then $y$ and
$\zeta\omega_{\cv\md,\tsl}$ lie on opposite sides of $\zeta/2$. Applied
coordinatewise to a vector $a \in \R^{\mathcal{C}}$, it gives
\[
  \norm{a-\zeta\omega}_{2}^{2} \;\ge\; \frac{\zeta^{2}}{4}\,d_{H}(\varpi(a),\omega) .
\]
For
(a), take $a_{\cv\md,\tsl} = \inner{\widehat{g}_{\cv\md}}{\wav{\lev}{\tsl}}$
and $\widehat\omega = \varpi(a)$: by Bessel's inequality,
$\lVert\widehat{g}_{\cv\md}-\gcomp{\cv}{\md}^{\omega}\rVert_{L_2[0,1]}^{2} \ge
\sum_{\tsl}(a_{\cv\md,\tsl}-\zeta\omega_{\cv\md,\tsl})^{2}$ for each
$(\cv,\md) \in \Kcal$, and it remains to sum over $\Kcal$. For (b), let
$\Pi_{\mathcal{C}}$ be the orthogonal projection of $L_2(P)$ onto the span of
the entries of $z_{\mathcal{C}}$, of dimension $|\mathcal{C}|$ because
$\Gram_{\mathcal{C}\mathcal{C}}$ is positive definite by (iii):
$\Pi_{\mathcal{C}}\widehat{f} = z_{\mathcal{C}}\T a$ with
$a = \Gram_{\mathcal{C}\mathcal{C}}^{-1}\E\{z_{\mathcal{C}}(\bX,\bU)\widehat{f}(\bX,\bU)\}$,
the expectation being over a new $(\bX,\bU)$ with $\widehat{f}$ fixed, so that
$a$ is a measurable function of the sample, $P$ being fixed. Since $f_{\omega}$
lies in the span, Pythagoras' theorem gives
$\norm{\widehat{f}-f_{\omega}}_{L_2(P)}^{2} \ge
\norm{\Pi_{\mathcal{C}}\widehat{f}-f_{\omega}}_{L_2(P)}^{2} =
(a-\zeta\omega)\T\Gram_{\mathcal{C}\mathcal{C}}(a-\zeta\omega) \ge
\gamma\norm{a-\zeta\omega}_{2}^{2}$; take $\widehat\omega = \varpi(a)$. For
(c), write $\widehat{\mathbf{F}} = (\widehat{f}(\bX_i,\bU_i))_{i\le n}$. On
$\mathcal{E}_{\mathcal{C}}$, $\bZ_{\mathcal{C}}$ has rank $|\mathcal{C}|$; take
$a = (\bZ_{\mathcal{C}}\T\bZ_{\mathcal{C}})^{-1}\bZ_{\mathcal{C}}\T\widehat{\mathbf{F}}$
and $\widehat\omega = \varpi(a)$ there, and $\widehat\omega = 0$ off it. On
$\mathcal{E}_{\mathcal{C}}$, since $\zeta\bZ_{\mathcal{C}}\omega$ is the vector
of the values of $f_{\omega}$ and lies in the column space of
$\bZ_{\mathcal{C}}$,
$n\normn{\widehat{f}-f_{\omega}}^{2} =
\norm{\widehat{\mathbf{F}}-\zeta\bZ_{\mathcal{C}}\omega}_{2}^{2} \ge
\norm{\bZ_{\mathcal{C}}(a-\zeta\omega)}_{2}^{2} =
n(a-\zeta\omega)\T\Gramhat_{\mathcal{C}\mathcal{C}}(a-\zeta\omega) \ge
n(\gamma/2)\norm{a-\zeta\omega}_{2}^{2}$; off it the right side of (c) is
zero. Finally $\Gramhat_{\mathcal{C}\mathcal{C}}$ is a principal submatrix of
$\Gramhat$ at resolution $J = \lev+1$, so that
$\lmin(\Gramhat_{\mathcal{C}\mathcal{C}}) \ge \lmin(\Gramhat)$, and
$\mathcal{E}_{\mathcal{C}}$ contains the event of Proposition~2(ii), which
gives the bound on its complement. \qed

\subsection*{Proof of Theorem~\ref{thm:lower}}

\emph{The hypercube.} Take $\zeta = \zeta_n$ and $\lev = \lev_n$ of
\eqref{eq:lower-choices}. Since $2^{\lev_n(2s+1)} \le C_g^{2}/\zeta_n^{2}$,
$\zeta_n2^{\lev_n(s+1/2)} \le C_g$, and Lemma~\ref{lem:hypercube}(i) gives
$f_{\omega} \in \mathcal{F}$ for every $\omega$. By
Lemma~\ref{lem:hypercube}(ii), the divergence between neighbors is at most one,
because $n\zeta_n^{2}B_X^{2}C_U = 2\sigma^{2}$, and
\eqref{eq:assouad} gives, for every $\widehat\omega$ with values in $\Omega$,
\begin{equation}
  \label{eq:assouad-n}
  \max_{\omega\in\Omega}\E_{\omega}d_{H}(\widehat\omega,\omega)
  \;\ge\; \cA|\mathcal{C}|,
  \qquad |\mathcal{C}| = |\Kcal|\,2^{\lev_n} .
\end{equation}

\emph{The scale.} By the maximality of $\lev_n$,
$2^{(\lev_n+1)(2s+1)} > C_g^{2}/\zeta_n^{2}$, so that
$2^{\lev_n} > (C_g/\zeta_n)^{2/(2s+1)}/2$ and
\begin{align*}
  2^{\lev_n}\zeta_n^{2} \;&>\; \frac12\,C_g^{2/(2s+1)}(\zeta_n^{2})^{2s/(2s+1)}
  = \frac12\,C_g^{2/(2s+1)}\Bigl(\frac{2\sigma^{2}}{nB_X^{2}C_U}\Bigr)^{2s/(2s+1)}
  \\
  &= 2^{-1/(2s+1)}\rlow^{2} \;\ge\; \frac12\,\rlow^{2},
\end{align*}
that is, $\zeta_n^{2}|\mathcal{C}| \ge |\Kcal|\rlow^{2}/2$.

\emph{Expectation.} The supremum over $\mathcal{F}$ is at least the maximum
over the hypercube. With the $\widehat\omega$ of
Lemma~\ref{lem:hypercube}(iv)(a) and \eqref{eq:assouad-n},
\[
  \max_{\omega}\E_{\omega}L_{\Kcal}(\widehat{g},f_{\omega}) \;\ge\;
  \frac{\zeta_n^{2}}{4}\,\cA|\mathcal{C}| \;\ge\; \frac{\cA}{8}\,|\Kcal|\,\rlow^{2} ;
\]
with that of (iv)(b), the same with the factor $\gamma$. For (iii), since
$d_{H} \le |\mathcal{C}|$ and the law of the $(\bX_i,\bU_i)$ does not depend
on $\omega$, $\E_{\omega}\{d_{H}(\widehat\omega,\omega)\mathbf{1}_{\mathcal{E}_{\mathcal{C}}}\}
\ge \E_{\omega}d_{H}(\widehat\omega,\omega) -
|\mathcal{C}|\Prob(\mathcal{E}_{\mathcal{C}}^{c})$, and for $n \ge n_{1}$
\[
  \max_{\omega}\E_{\omega}\normn{\widehat{f}-f_{\omega}}^{2}
  \;\ge\; \frac{\gamma\zeta_n^{2}}{8}\Bigl(\cA|\mathcal{C}| -
  \frac{\cA}{4}|\mathcal{C}|\Bigr)
  = \frac{3\cA}{32}\,\gamma\zeta_n^{2}|\mathcal{C}|
  \;\ge\; \frac{\cA}{32}\,\gamma|\Kcal|\,\rlow^{2} .
\]

\emph{Probability.} Let $\omega^{*}$ attain the maximum in
\eqref{eq:assouad-n}, which exists because $\Omega$ is finite. Since
$0 \le d_{H} \le |\mathcal{C}|$,
$\cA|\mathcal{C}| \le \E_{\omega^{*}}d_{H} \le \cA|\mathcal{C}|/2 +
|\mathcal{C}|\Prob_{\omega^{*}}(d_{H} \ge \cA|\mathcal{C}|/2)$, so that
$\Prob_{\omega^{*}}(d_{H} \ge \cA|\mathcal{C}|/2) \ge \cA/2$. On that event,
(iv)(a) gives $L_{\Kcal} \ge (\zeta_n^{2}/4)\cA|\mathcal{C}|/2 \ge
(\cA/16)|\Kcal|\rlow^{2}$, and (iv)(b) the same with the factor $\gamma$; on
its intersection with $\mathcal{E}_{\mathcal{C}}$, of probability at least
$\cA/2 - \cA/4$ when $n \ge n_{1}$, (iv)(c) gives
$\normn{\widehat{f}-f_{\omega^{*}}}^{2} \ge
(\gamma\zeta_n^{2}/8)\cA|\mathcal{C}|/2 \ge (\cA/32)\gamma|\Kcal|\rlow^{2}$.
The estimator being arbitrary, the infima follow.

\emph{The integer $n_{1}$ is finite.} Since $2^{\lev_n+1} \le
2(C_g/\zeta_n)^{2/(2s+1)}$, of order $n^{1/(2s+1)}$, $D$ and $R_J$ with
$J = \lev_n + 1$ are of order $n^{1/(2s+1)}$, and the bound of
Lemma~\ref{lem:hypercube}(iv)(c) is of order
$n^{1/(2s+1)}\exp(-cn^{2s/(2s+1)})$ for a constant $c > 0$, which tends to
zero. The two particular cases are $\Kcal$ a single block and $\Kcal$ all the
blocks. \qed

% E5f (E1.15; D54): derivations/08-blocos.tex, Corolario 15 (numeracao global;
% cor:truncada-blocos) e a sua prova, dentro do grupo azul de k = 4.
\subsection*{The slow rate through a truncated comparator}

Proposition~3 reads Theorem~2(i) at the comparator $\betastar$ only, and pays
$\pen w_{G}\norm{\thetastar_{G}}_{2}$ for every chunk, including the chunks
whose energy would cost less as bias. Theorem~2 holds for every comparator on
the same event, and the comparator truncated chunk by chunk,
\begin{equation}
  \label{eq:truncated-blocks}
  \thbar_{G} \;=\; \thetastar_{G}\,\mathbf{1}\bigl\{\norm{\thetastar_{G}}_{2}
  > \pen w_{G}\bigr\},
  \qquad G \in \Gcal,
\end{equation}
charges the penalty only for the chunks above their own level and sends the
others to the bias. What remains is
\begin{equation}
  \label{eq:RG}
  \RGw(\btheta;\pen) = \sum_{G\in\Gcal}
  \min\bigl(\pen w_{G}\norm{\btheta_{G}}_{2},\,\norm{\btheta_{G}}_{2}^{2}\bigr),
  \qquad
  \RGone(\btheta;t) = \sum_{G\in\Gcal}
  \min\bigl(t\norm{\btheta_{G}}_{2},\,\norm{\btheta_{G}}_{2}^{2}\bigr),
\end{equation}
the counterparts in the group norm of the ideal risk $\Rid$, with
$\RGone(\btheta;t) \le t\nGone{\btheta}$.

\begin{proposition}[Slow rate through the truncated comparator]
\label{prop:truncated}
Let the assumptions of Proposition~3 hold; in the statements with $\lamGp$
in the estimator, the conditions on $J_n$ are replaced by $2^{J_n} \le n^{2}$,
as in Proposition~\ref{prop:slow-ceiling}.
\begin{enumerate}
\item[(i)] On the event $\Tw$, with $\thbar$ of \eqref{eq:truncated-blocks},
\[
  \normn{\widehat{f}-f}^{2}
  \;\le\; 18\pen\nGw{\thbar} + 48\normn{\Bmat(\thetastar-\thbar)}^{2}
  + 48\normn{\bb}^{2} + 3\normn{\PA\bveps}^{2} .
\]
If $\pen$ and $w$ are deterministic,
$\E\normn{\Bmat(\thetastar-\thbar)}^{2} \le
\kappa_2C_U\norm{\thetastar-\thbar}_{2}^{2}$,
\[
  18\pen\nGw{\thbar} + 48\kappa_2C_U\norm{\thetastar-\thbar}_{2}^{2}
  \;\le\; \max(18,48\kappa_2C_U)\,\RGw(\thetastar;\pen),
\]
and $\RGw(\thetastar;\lamG) \le \bar\varrho\,\RGone(\thetastar;\lamGone)$.
Consequently, with deterministic weights,
\[
  \normn{\widehat{f}-f}^{2} \;=\; O_p\bigl\{\RGone(\thetastar;\lamGone)
  + \mathcal{B}_{J_n}^{2} + p/n\bigr\},
\]
and, with $\lamGp$ in the estimator, the same holds with
$(\mu^{\Gcal}_{J_n})^{1/2}\lamGone$ in place of $\lamGone$, for every sequence
$J_n$.
\item[(ii)] For every $J$, every $b \ge 2$ and every $t > 0$, with
$\tau = (s+1/2)^{-1}$: if $s < 1/2$,
\begin{equation}
  \label{eq:count-truncated}
  \RGone(\thetastar;t) \;\le\; pq\Bigl\{\frac{C_g\,t}{1-2^{-\seff}}
  + C_g^{\tau}\,t^{4s/(2s+1)}\,b^{\{(2/\pi-1)_{+}-2s\}/(2s+1)}
  \Bigl(\frac{1}{1-2^{s-1/2}} + \frac{1}{1-2^{-\min(\pi,2)\seff}}\Bigr)\Bigr\};
\end{equation}
if $s \ge 1/2$, nothing is gained in order: $\RGone(\thetastar;t) \le
t\nGone{\thetastar}$, and for $t \le C_g$ the class of Assumption~5 contains a
$\thetastar$ with $\RGone(\thetastar;t) \ge tC_g$ if $s > 1/2$, and, if
$s = 1/2$, one in which every fine level $\lev < J$ with
$2^{\lev} \le C_gb^{1-\vartheta_{\pi}}/t$ contributes at least
$tC_gb^{-\vartheta_{\pi}}/2$ to $\RGone(\thetastar;t)$, the order of
Lemma~\ref{lem:count}.
\item[(iii)] If $s < 1/2$ and $\seff > s/(2s+1)$, then with
$J_n = \lceil c\log_2 n\rceil$ and $c$ in the window (3.2) of Theorem~1,
\[
  \normn{\widehat{f}-f}^{2}
  \;=\; O_p\bigl\{n^{-2s/(2s+1)}(\log n)^{(2/\pi-1)_{+}/(2s+1)}\bigr\},
\]
the prediction rate of Theorem~1.
\item[(iv)] If $s < 1/2$ and $\seff \le s/(2s+1)$, which forces $\pi < 2$,
then with $\lamGp$ in the estimator and the choice
$2^{J_n} \asymp \{n^{2s}\bn^{-(2s+2/\pi-1)/2}\}^{1/\{s+\seff(2s+1)\}}$, which
exceeds $n/\bn$,
\[
  \normn{\widehat{f}-f}^{2}
  \;=\; O_p\bigl\{n^{-e}(\log n)^{e(2s+2/\pi-1)/(4s)}\bigr\},
  \qquad e = \frac{4s\seff}{s+\seff(2s+1)} ;
\]
at $\seff = s/(2s+1)$ this is
$n^{-2s/(2s+1)}(\log n)^{(2s+2/\pi-1)/\{2(2s+1)\}}$.
\end{enumerate}
\end{proposition}

Proposition~\ref{prop:truncated} is a statement on prediction. Like
Proposition~3, it uses no lower bound on $\Gramhat$: besides the upper bound
on $\lmaxx(\Gramhat)$ that makes the level admissible, it uses only
$\lmaxx(\Gram) \le \kappa_2C_U$, in the new bias, and the components still
need the design condition (Corollary~1). When $s < 1/2$ and $\pi < 2$ it
improves on Proposition~3(ii) and on Proposition~\ref{prop:slow-ceiling}(iii)(b)
in every pair: (iii) gives the rate of Theorem~1 in the window of that
theorem, and (iv) replaces Proposition~\ref{prop:slow-ceiling}(iii)(b) and the
part of Proposition~3(ii) with $s/2 < \seff \le s/(2s+1)$. When $\pi \ge 2$
the rate is that of Proposition~3(ii), $n^{-2s/(2s+1)}$, now for every
$c \in [1/(2s+1),1)$. When $s \ge 1/2$ nothing changes: a coarse chunk of
norm $C_g$, which the class allows, costs $\pen C_g$ in the penalty or $C_g^{2}$ in the bias, and
the slow route does not go below $\lamGone \asymp (\log n/n)^{1/2}$;
Proposition~3(i) and Proposition~\ref{prop:slow-ceiling}(iii)(a) and (iv)
remain the statements with no design condition in those pairs. Against the
lasso read at its own truncated comparator
(Proposition~\ref{prop:lasso-truncated}), the chunks remove the factor
$(\log n)^{2\seff/(2s+1)}$ in (iii), the gain of Theorem~1 over
Theorem~\ref{thm:lasso-main}, and $(\log n)^{e\seff/(2s)}$ in (iv). The
proposition is stated for the estimator (2.5) and not for the one that the
software computes (Remark~\ref{rem:white}); for the latter the comparator
would truncate at $\pen w_{G}\Lambda^{1/2}$ on the event $q_{\max} \le \Lambda$,
and we do not pursue it.

\subsection*{Proof of Proposition~\ref{prop:truncated}}

\emph{Part (i).} Theorem~2(i) and (iii) with $\bar\bbeta = (\bc,\thbar)$ give,
as in the proof of \eqref{eq:slow-nonasymptotic},
$\normn{\widehat{f}-f}^{2} \le 18\pen\nGw{\thbar} + 24\normn{\bbar}^{2} +
3\normn{\PA\bveps}^{2}$, and $\bbar = f - \Amat\bc - \Bmat\thbar = \bb +
\Bmat(\thetastar-\thbar)$ gives $\normn{\bbar}^{2} \le 2\normn{\bb}^{2} +
2\normn{\Bmat(\thetastar-\thbar)}^{2}$. If $\pen$ and $w$ are deterministic,
so is $u = \thetastar-\thbar$, and $\E\normn{\Bmat u}^{2} =
u\T\E(\Bmat\T\Bmat/n)u \le \lmaxx(\Gram)\norm{u}_{2}^{2} \le
\kappa_2C_U\norm{u}_{2}^{2}$, because $\E(\Bmat\T\Bmat/n)$ is a principal
submatrix of $\Gram$, by the upper half of Proposition~2(i). On the chunks with
$\norm{\thetastar_{G}}_{2} > \pen w_{G}$ the minimum in \eqref{eq:RG} is
$\pen w_{G}\norm{\thetastar_{G}}_{2}$, and on the others it is
$\norm{\thetastar_{G}}_{2}^{2}$; summing gives the bound by
$\max(18,48\kappa_2C_U)\RGw(\thetastar;\pen)$. By
Lemma~\ref{lem:calib-blocks}(iii), $\lamG w_{G} \le \varrho\lamGone$, and
$\min(\varrho x,y) \le \varrho\min(x,y)$ for $\varrho \ge 1$ gives
$\RGw(\thetastar;\lamG) \le \varrho\,\RGone(\thetastar;\lamGone)$; with
$\lamGp$, the same with $(\mu^{\Gcal}_{J})^{1/2}\lamGone$. The form in $O_p$
is obtained as in the proof of Proposition~3, on the intersection of
$\Acal^{+}$ below the ceiling, or of $\Acal^{++}$ of the proof of
Proposition~\ref{prop:slow-ceiling} for every $J_n$, with $\Ecal$ and with
the events given by Markov's inequality for $\normn{\bb}^{2}$,
$\normn{\Bmat(\thetastar-\thbar)}^{2}$ and $\normn{\PA\bveps}^{2}$.

\emph{Part (ii), $s < 1/2$.} Fix a block. The coarse chunk contributes at most
$t\norm{\theta_{\mathrm{coarse}}}_{2} \le tC_g(1-2^{-\seff})^{-1}$, by the
proof of Lemma~\ref{lem:count}. At a fine level $\lev$, with $\Gcal_{\lev}$
its chunks and $v_{\lev} = \sum_{G\in\Gcal_{\lev}}\min(t\norm{\theta_{G}}_{2},
\norm{\theta_{G}}_{2}^{2})$, there are two bounds. The first is in the group
norm: $v_{\lev} \le t\sum_{G}\norm{\theta_{G}}_{2} \le
tC_gb^{-\vartheta_{\pi}}2^{\lev(1/2-s)} =: A_{\lev}$, by the proof of
Lemma~\ref{lem:count}. The second is the tail. If $\pi \le 2$,
$\min(ty,y^{2}) \le t^{2-\pi}y^{\pi}$ for $y \ge 0$ (if $y \le t$,
$y^{2} = y^{\pi}y^{2-\pi} \le y^{\pi}t^{2-\pi}$; if $y > t$,
$ty = t^{2-\pi}t^{\pi-1}y \le t^{2-\pi}y^{\pi}$, because $\pi \ge 1$), and
$\norm{\theta_{G}}_{2} \le \norm{\theta_{G}}_{\pi}$ gives
$v_{\lev} \le t^{2-\pi}\sum_{G}\norm{\theta_{G}}_{\pi}^{\pi} =
t^{2-\pi}\norm{\theta_{\lev}}_{\pi}^{\pi} \le
t^{2-\pi}C_g^{\pi}2^{-\lev\pi\seff}$, since $\pi(s+1/2-1/\pi) = \pi\seff$; if
$\pi \ge 2$, $v_{\lev} \le \norm{\theta_{\lev}}_{2}^{2} \le C_g^{2}2^{-2\lev s}$.
Call this second bound $B_{\lev}$. With
\[
  2^{x} = \bigl\{(C_g/t)\,b^{1-\vartheta_{\pi}}\bigr\}^{\tau},
  \qquad
  h = C_g^{\tau}\,t^{2-\tau}\,b^{\tau(1-\vartheta_{\pi})-1},
\]
one has $A_{\lev} = h2^{(\lev-x)(1/2-s)}$ and
$B_{\lev} = h2^{-(\lev-x)\min(\pi,2)\seff}$. Indeed, at $\lev = x$,
$A_{x} = C_g^{\tau}t^{2-\tau}b^{-\vartheta_{\pi}+(1-\vartheta_{\pi})(\tau-1)}$,
because $\tau(1/2-s) = \tau-1$, and the exponent of $b$ is
$\tau(1-\vartheta_{\pi})-1$; if $\pi \le 2$, $1-\vartheta_{\pi} = 1/\pi$ and
$B_{x} = t^{2-\pi}C_g^{\pi}\{(C_g/t)b^{1/\pi}\}^{-(\pi-\tau)} =
C_g^{\tau}t^{2-\tau}b^{\tau/\pi-1}$, because $\tau\pi\seff = \pi-\tau$; if
$\pi \ge 2$, $1-\vartheta_{\pi} = 1/2$ and
$B_{x} = C_g^{2}\{(C_g/t)b^{1/2}\}^{-(2-\tau)} =
C_g^{\tau}t^{2-\tau}b^{\tau/2-1}$, because $2s\tau = 2-\tau$. For every $\pi$,
$\tau(1-\vartheta_{\pi})-1 = \{(2/\pi-1)_{+}-2s\}/(2s+1)$ and
$2-\tau = 4s/(2s+1)$. Summing $A_{\lev}$ over the fine levels with
$\lev \le x$ and $B_{\lev}$ over those with $\lev > x$, two geometric series
of ratios $2^{-(1/2-s)}$ and $2^{-\min(\pi,2)\seff}$ that start from at most
$h$, gives at most $h\{(1-2^{s-1/2})^{-1} + (1-2^{-\min(\pi,2)\seff})^{-1}\}$
per block. Truncating at $\lev < J$ only removes terms, and summing over the
$pq$ blocks gives \eqref{eq:count-truncated}.

\emph{Part (ii), $s \ge 1/2$.} The first statement is
$\min(ty,y^{2}) \le ty$. If $s > 1/2$, the block with
$\theta^{*}_{\cv\md,00} = C_g$ and all its other coefficients zero satisfies
Assumption~5, and its coarse chunk, of norm $C_g$, contributes
$\min(tC_g,C_g^{2}) = tC_g$. If $s = 1/2$, take a fine level $\lev$ with
$2^{\lev} \le C_gb^{1-\vartheta_{\pi}}/t$, and let $k_{\lev}$ be the number of
its chunks, so that $2^{\lev}/(2b) \le k_{\lev} \le 2^{\lev}/b$. If
$\pi \le 2$, put a spike of height
$y_{\lev} = C_g2^{-\lev(1-1/\pi)}k_{\lev}^{-1/\pi}$ in each chunk, which
saturates Assumption~5 at that level; since $k_{\lev} \le 2^{\lev}/b$,
$y_{\lev} \ge C_g2^{-\lev}b^{1/\pi} \ge t$, and the level contributes
$k_{\lev}ty_{\lev} = tC_g2^{-\lev(1-1/\pi)}k_{\lev}^{1-1/\pi} \ge
tC_g(2b)^{-(1-1/\pi)}$. If $\pi \ge 2$, set all the coefficients of the level
equal to $C_g2^{-\lev}$, which again saturates Assumption~5: each chunk has
norm $|G|^{1/2}C_g2^{-\lev} \ge b^{1/2}C_g2^{-\lev} \ge t$, and the level
contributes $tC_g2^{-\lev}\sum_{G}|G|^{1/2} \ge tC_g2^{-\lev}k_{\lev}b^{1/2}
\ge tC_gb^{-1/2}/2$. In both cases the contribution is at least
$tC_gb^{-\vartheta_{\pi}}/2$.

\emph{Part (iii).} With $\pen = \lamG$ and
$(\lamGone)^{2} \le C_{\pen}\sigma^{2}\bn/n$ by \eqref{eq:order-lambda},
\eqref{eq:count-truncated} with $b = \bn$ gives
$\RGone(\thetastar;\lamGone) \lesssim (\bn/n)^{1/2} +
(\bn/n)^{2s/(2s+1)}\bn^{\{(2/\pi-1)_{+}-2s\}/(2s+1)} = (\bn/n)^{1/2} +
n^{-2s/(2s+1)}\bn^{(2/\pi-1)_{+}/(2s+1)}$, and the first term is of smaller
order because $2s/(2s+1) < 1/2$. The bias and the conditions on $J_n$ are
those of the proof of Theorem~1: $c \ge s/\{(2s+1)\seff\}$ gives
$\mathcal{B}_{J_n}^{2} \lesssim n^{-2s/(2s+1)}$, and $c < 1$ gives
$2^{J_n} \le 2n$ and $pq2^{J_n}\log(pq2^{J_n})/n \to 0$.

\emph{Part (iv).} Write $\eta = 2s/\{s+\seff(2s+1)\} \in [1,2)$ and
$\beta = (2s+2/\pi-1)/(2s+1)$, which lies in $(0,\,4s/(2s+1))$ because
$\pi < 2$ and $2\seff = 2s+1-2/\pi > 0$. The choice of the statement is
$2^{J_n} \asymp n^{\eta}\bn^{-\eta\beta(2s+1)/(4s)}$, which exceeds $n/\bn$
because $\beta(2s+1)/(4s) < 1$, and is at most $n^{2}$. By
\eqref{eq:lam-plus-blocks}, $t = (\mu^{\Gcal}_{J_n})^{1/2}\lamGone \asymp
2^{J_n/2}\bn/n$, and \eqref{eq:count-truncated} gives
\[
  t^{4s/(2s+1)}\bn^{\beta-4s/(2s+1)}
  \;\asymp\; 2^{J_n\cdot 2s/(2s+1)}\,n^{-4s/(2s+1)}\,\bn^{\beta},
\]
which, equated with the bias $2^{-2J_n\seff}$, gives
$2^{J_n\{2s/(2s+1)+2\seff\}} = n^{4s/(2s+1)}\bn^{-\beta}$; this is the choice
of the statement, because $\{2s/(2s+1)+2\seff\}(2s+1) = 2\{s+\seff(2s+1)\}$.
The common value is $2^{-2J_n\seff} =
\{\bn^{\beta}/n^{4s/(2s+1)}\}^{2\seff/\{2s/(2s+1)+2\seff\}} =
n^{-e}\bn^{e\beta(2s+1)/(4s)}$, and $\beta(2s+1)/(4s) = (2s+2/\pi-1)/(4s)$.
The term of the coarse chunk, of order $t$, is polynomially smaller, because
$t \to 0$ and $4s/(2s+1) < 1$. \qed


%%%%%%%%%%%%%%%%%%%%%%%%%%%%%%%%%%%%%%%%%%%%%%%%%%%%%%%%%%%%%%%%%%%%%%%%%%%%%%
% derivations/06-selecao-limiar.tex (Lema 13 e Corolario 8 na numeracao
% global; a Hipotese S e a Assumption 6 do artigo). Entrou em k = 2 (E5c).

